\begin{examplebox}
	\textbf{\textcolor{CExample}{例 1（基础）}}
	\textbf{\textcolor{CExample}{题目：}}
	已知圆$O$半径$OA=3$，点$P$在圆外，$OP=5$，$PA$、$PB$是圆$O$的两条切线，切点分别为$A$、$B$。求$PA$、$PB$的长度。
	\textbf{\textcolor{CExample}{解题步骤：}}
	\begin{description}[style=nextline, leftmargin=2.8em, labelsep=0.5em, itemsep=0.45em]
		\item[\textbf{第一步（垂直条件）：}]
			连接$OA$，由切线性质，$OA\perp PA$，故$\triangle OAP$为直角三角形。
			\[
				\angle OAP = 90^\circ.
			\]
			\par\textbf{理由：} 圆的切线垂直于过切点的半径。
		\item[\textbf{第二步（勾股计算）：}]
			在$\mathrm{Rt}\triangle OAP$中，利用勾股定理：
			\[
				\begin{aligned}
					PA &= \sqrt{OP^2 - OA^2} \\ &= \sqrt{5^2 - 3^2} \\ &= \sqrt{25-9} \\ &= \sqrt{16} \\ &= 4.
			\end{aligned}
			\]
			\par\textbf{理由：} 直角三角形中斜边平方等于两直角边平方和。
		\item[\textbf{第三步（使用定理）：}]
			由切线长定理，$PA = PB$，所以$PB = 4$。
			\[
				PB = PA = 4.
			\]
			\par\textbf{理由：} 切线长定理：从圆外一点引两条切线，切线长相等。
	\end{description}
	\textbf{\textcolor{CExample}{关键结论：}}
	\textbf{$PA = PB = 4$。}

	\medskip
	\textbf{\textcolor{CExample}{例 2（稍变形）}}
	\textbf{\textcolor{CExample}{题目：}}
	如图，$PA$、$PB$是$\odot O$的两条切线，$A$、$B$为切点，$\angle APB = 60^\circ$，$\odot O$半径$OA=2$。求$PA$的长。
	\textbf{\textcolor{CExample}{解题步骤：}}
	\begin{description}[style=nextline, leftmargin=2.8em, labelsep=0.5em, itemsep=0.45em]
		\item[\textbf{第一步（角平分）：}]
			由切线长定理，$OP$平分$\angle APB$，故$\angle APO = \frac12 \angle APB = 30^\circ$。
			\[
				\angle APO = 30^\circ.
			\]
			\par\textbf{理由：} 切线长定理：圆心与圆外点的连线平分两条切线所夹的角。
		\item[\textbf{第二步（三角函数）：}]
			在$\mathrm{Rt}\triangle OAP$中，$\angle OAP = 90^\circ$，$\angle APO = 30^\circ$，$OA=2$。利用$\tan\angle APO = \frac{OA}{PA}$，得
			\[
				\begin{aligned}
					PA &= \frac{OA}{\tan 30^\circ} \\ &= \frac{2}{\frac{\sqrt{3}}{3}} \\ &= 2\sqrt{3}.
			\end{aligned}
			\]
			\par\textbf{理由：} 直角三角形中，正切等于对边比邻边。
		\item[\textbf{第三步（切线长相等）：}]
			由切线长定理，$PA = PB$，故$PA = 2\sqrt{3}$。
			\[
				PA = 2\sqrt{3}.
			\]
			\par\textbf{理由：} 切线长定理：从圆外一点引两条切线，切线长相等。
	\end{description}
	\textbf{\textcolor{CExample}{关键结论：}}
	\textbf{$PA = 2\sqrt{3}$。}

	\medskip
	\textbf{\textcolor{CExample}{例 3（综合应用）}}
	\textbf{\textcolor{CExample}{题目：}}
	已知$PA$、$PB$是$\odot O$的切线，$A$、$B$是切点，$OP$与$AB$交于点$C$。若$PA=5$，$OA=3$，求$AB$的长。
	\textbf{\textcolor{CExample}{解题步骤：}}
	\begin{description}[style=nextline, leftmargin=2.8em, labelsep=0.5em, itemsep=0.45em]
		\item[\textbf{第一步（垂直平分）：}]
			由切线长定理的推广，$OP$垂直平分$AB$，故$AC=BC=\frac12 AB$，且$OP\perp AB$。
			\[
				AC = BC,\quad OP\perp AB.
			\]
			\par\textbf{理由：} 由$OA=OB$、$PA=PB$可证$OP$是$AB$的垂直平分线（全等或两点确定垂直平分线）。
		\item[\textbf{第二步（求OP）：}]
			在$\mathrm{Rt}\triangle OAP$中，$OA\perp PA$，由勾股定理：
			\[
				\begin{aligned}
					OP &= \sqrt{OA^2 + PA^2} \\ &= \sqrt{3^2 + 5^2} \\ &= \sqrt{34}.
			\end{aligned}
			\]
			\par\textbf{理由：} 直角三角形两直角边平方和等于斜边平方。
		\item[\textbf{第三步（面积法求AB）：}]
			利用$\triangle OAP$的面积公式：$S = \frac12 OA\cdot PA = \frac12 OP\cdot AC$，所以
			\[
				\begin{aligned}
					AC &= \frac{OA\cdot PA}{OP} \\ &= \frac{3\times5}{\sqrt{34}} \\ &= \frac{15}{\sqrt{34}},\\ AB &= 2AC \\ &= \frac{30}{\sqrt{34}} \\ &= \frac{15\sqrt{34}}{17}.
			\end{aligned}
			\]
			\par\textbf{理由：} 直角三角形面积可以用两种直角边乘积的一半，也可以用斜边与斜边上高的乘积的一半。
	\end{description}
	\textbf{\textcolor{CExample}{关键结论：}}
	\textbf{$AB = \dfrac{15\sqrt{34}}{17}$。}
\end{examplebox}
